\section{Objectifs}
\begin{itemize}
\item Reconnaître une modélisation quadratique.
\item Lire les éléments caractéristiques d'une parabole.
\item Exploiter une forme factorisée.
\item Résoudre un problème simple de maximum, minimum ou signe.
\end{itemize}

\section{1. Fonctions quadratiques}
Une fonction polynôme de degré 2 s'écrit
\[
f(x)=ax^2+bx+c,\qquad a\ne0.
\]

Sa courbe est une \textbf{parabole}.

Le programme conduit notamment à manipuler des formes simples comme
\[
ax^2,\qquad ax^2+c,\qquad a(x-x_1)(x-x_2).
\]

\section{2. Axe de symétrie et sommet}
Une parabole possède un axe de symétrie vertical. Le sommet est le point où la fonction atteint un minimum si la parabole est tournée vers le haut, ou un maximum si elle est tournée vers le bas.

Aucune formule générale du sommet n'est exigée ici. On peut le déterminer par lecture graphique, symétrie ou raisonnement adapté à l'expression.

\section{3. Racines et forme factorisée}
Si
\[
f(x)=a(x-x_1)(x-x_2),
\]
alors
\[
f(x_1)=0\qquad\text{et}\qquad f(x_2)=0.
\]

Les nombres $x_1$ et $x_2$ sont les racines.

Le calcul des racines à l'aide du discriminant n'est pas au programme de ce module.

\section{4. Signe d'une forme factorisée}
Pour
\[
f(x)=(x-2)(x-5),
\]
on obtient :
\begin{itemize}
\item $f(x)>0$ pour $x<2$ ;
\item $f(2)=0$ ;
\item $f(x)<0$ pour $2<x<5$ ;
\item $f(5)=0$ ;
\item $f(x)>0$ pour $x>5$.
\end{itemize}

On justifie ce résultat à partir du signe de chaque facteur.

\section{5. Modéliser une situation}
Une parabole peut servir à représenter une trajectoire simplifiée, une arche ou une grandeur qui augmente puis diminue. Comme tout modèle, elle n'est pertinente que sur un domaine compatible avec la situation réelle.

\section{Automatismes à entretenir}
\begin{itemize}
\item Développer, réduire et factoriser une expression simple.
\item Résoudre une équation produit nul.
\item Lire le signe et les variations sur un graphique.
\item Résoudre graphiquement $f(x)=k$.
\end{itemize}

\section{Erreurs fréquentes}
\begin{itemize}
\item Employer le discriminant alors qu'il n'est pas attendu ici.
\item Confondre racine d'une fonction et racine carrée.
\item Oublier de limiter l'interprétation au domaine réel du problème.
\end{itemize}

\section{À retenir}
La modélisation quadratique fournit une première famille de modèles non linéaires simples, adaptée aux phénomènes présentant un maximum ou un minimum.
